\documentclass[12pt]{article}

\usepackage[utf8]{inputenc}
\usepackage[T1]{fontenc}
\usepackage{mathptmx}
\usepackage[margin=1in]{geometry}
\usepackage{booktabs}
\usepackage{longtable}
\usepackage{array}
\usepackage{caption}

\captionsetup{font=small, labelfont=bf}
\setlength{\parindent}{0pt}
\newcolumntype{R}[1]{>{\raggedleft\arraybackslash}p{#1}}

\title{}
\date{}

\begin{document}
\pagestyle{empty}
\setcounter{table}{1}

\begin{longtable}{p{6.0cm}R{3.5cm}R{3.5cm}}
\caption{Included versus excluded participants, among the 3,124 outcome-eligible participants.}
\label{tab:t1b}\\
\toprule
Characteristic & Included ($n=1{,}784$) & Excluded ($n=1{,}340$) \\
\midrule
\endfirsthead
\toprule
Characteristic & Included ($n=1{,}784$) & Excluded ($n=1{,}340$) \\
\midrule
\endhead
\bottomrule
\endlastfoot
Age, years, mean (SD) & 69.37 (6.74) & 70.10 (6.94) \\
Income-to-poverty ratio, mean (SD) & 2.60 (1.60) & 2.50 (1.59) \\
PHQ-9 depression score, median [IQR] & 1.00 [0.00, 4.00] & 2.00 [0.00, 5.00] \\
Female, \% & 50.7\% & 52.6\% \\
Low cognitive performance, \% & 24.2\% & 26.1\% \\
Diabetes, \% & 27.5\% & 30.0\% \\
\end{longtable}

\end{document}
